\chapter{Straight Lines in Space and Skew Lines (Pages 37--52)}
\label{ch:straight_line_skew}

% =========================================================================
% PAGE 37
% =========================================================================
\section{Page 37: The Straight Line in Space --- Symmetrical and Unsymmetrical Forms}
\label{sec:page37}

\begin{conceptbox}[Manuscript Overview: Page 37]
Page 37 introduces the straight line in three-dimensional space. It defines the unsymmetrical form (intersection of two non-parallel planes) and derives the symmetrical (point-direction) form and two-point form of a straight line.
\end{conceptbox}

\subsection{The Line as the Intersection of Two Planes}
In three-dimensional analytical geometry, a single linear equation $Ax + By + Cz + D = 0$ represents a plane. A straight line is the locus of points common to two intersecting planes:
\begin{equation}
    \begin{cases}
    A_1 x + B_1 y + C_1 z + D_1 = 0, \\
    A_2 x + B_2 y + C_2 z + D_2 = 0,
    \end{cases} \quad \text{provided } \vec{N}_1 \nparallel \vec{N}_2.
\end{equation}
This representation is known as the \textbf{unsymmetrical form} (or general form) of a straight line.

\subsection{Derivation of the Symmetrical Form}
\begin{theorembox}[Symmetrical Form of a Straight Line]
The equations of a straight line passing through a given point $P_1(x_1, y_1, z_1)$ and having direction cosines $(l, m, n)$ are:
\begin{equation}
    \boxed{\frac{x - x_1}{l} = \frac{y - y_1}{m} = \frac{z - z_1}{n} = r}
\end{equation}
where $r$ represents the signed algebraic distance along the line from $P_1$ to any point $P(x, y, z)$.
\end{theorembox}

\begin{proof}
Let $P_1(x_1, y_1, z_1)$ be a fixed point on the line, and let $P(x, y, z)$ be any arbitrary point on the line.
Let $r$ denote the directed distance from $P_1$ to $P$, so that $P_1 P = r$.
Let $(l, m, n)$ be the direction cosines of the line, which means:
\begin{equation}
    l = \cos\alpha, \quad m = \cos\beta, \quad n = \cos\gamma.
\end{equation}
The projections of the segment $P_1 P$ on the coordinate axes are (Page 17):
\begin{align}
    x - x_1 &= r\cos\alpha = lr, \\
    y - y_1 &= r\cos\beta = mr, \\
    z - z_1 &= r\cos\gamma = nr.
\end{align}
Solving each equation for $r$:
\begin{equation}
    r = \frac{x - x_1}{l} = \frac{y - y_1}{m} = \frac{z - z_1}{n}.
\end{equation}
Since $r$ is identical for all three projections:
\begin{equation}
    \boxed{\frac{x - x_1}{l} = \frac{y - y_1}{m} = \frac{z - z_1}{n} = r}
\end{equation}
If direction ratios $(a, b, c)$ are given instead of direction cosines:
\begin{equation}
    \boxed{\frac{x - x_1}{a} = \frac{y - y_1}{b} = \frac{z - z_1}{c} = \lambda}
\end{equation}
where $\lambda = \frac{r}{\sqrt{a^2+b^2+c^2}}$.
\end{proof}

\subsection{Two-Point Form of a Straight Line}
If a straight line passes through two given points $P_1(x_1, y_1, z_1)$ and $P_2(x_2, y_2, z_2)$, its direction ratios are $(x_2 - x_1, \; y_2 - y_1, \; z_2 - z_1)$ (Page 21). Substituting these ratios yields:
\begin{equation}
    \boxed{\frac{x - x_1}{x_2 - x_1} = \frac{y - y_1}{y_2 - y_1} = \frac{z - z_1}{z_2 - z_1}}
\end{equation}

% =========================================================================
% PAGE 38
% =========================================================================
\section{Page 38: Parametric Coordinates and the General Point on a Line}
\label{sec:page38}

\begin{conceptbox}[Manuscript Overview: Page 38]
Page 38 formalizes the parametric representation of a line and details how coordinates of any general point on a line can be expressed in terms of a single scalar parameter.
\end{conceptbox}

\subsection{Parametric Coordinates}
From the symmetrical form $\frac{x - x_1}{l} = \frac{y - y_1}{m} = \frac{z - z_1}{n} = r$:
\begin{align}
    x &= x_1 + lr, \nonumber\\
    y &= y_1 + mr, \nonumber\\
    z &= z_1 + nr. \label{eq:param_coords}
\end{align}
Equations \eqref{eq:param_coords} are the \textbf{parametric equations} of the line.

\begin{theorembox}[General Point on a Line]
The coordinates of any point $P$ on the line passing through $(x_1, y_1, z_1)$ with direction ratios $(a, b, c)$ can be expressed in terms of a real parameter $\lambda$ as:
\begin{equation}
    \boxed{P(\lambda) = (x_1 + a\lambda, \; y_1 + b\lambda, \; z_1 + c\lambda), \quad \lambda \in \mathbb{R}}
\end{equation}
If $(a, b, c)$ are normalized direction cosines $(l, m, n)$, then $|\lambda| = r$ is the exact Euclidean distance from $(x_1, y_1, z_1)$ to $P$.
\end{theorembox}

\subsection{Geometrical Significance}
This parametric form reduces three-dimensional geometric problems concerning lines (such as intersections, perpendicular feet, reflections, and distances) to single-variable algebraic equations in $\lambda$.

% =========================================================================
% PAGE 39
% =========================================================================
\section{Page 39: Transformation from Unsymmetrical Form to Symmetrical Form}
\label{sec:page39}

\begin{conceptbox}[Manuscript Overview: Page 39]
Page 39 provides the complete mathematical algorithm for transforming a straight line given in unsymmetrical form into the standard symmetrical form.
\end{conceptbox}

\subsection{The Reduction Algorithm}
Let the line be given as the intersection of two non-parallel planes:
\begin{align}
    \Pi_1 &: A_1 x + B_1 y + C_1 z + D_1 = 0, \label{eq:p39_pi1} \\
    \Pi_2 &: A_2 x + B_2 y + C_2 z + D_2 = 0. \label{eq:p39_pi2}
\end{align}

\subsubsection{Step 1: Finding the Direction Ratios of the Line}
Let $(l, m, n)$ be the direction ratios of the line of intersection $L$.
Since $L$ lies completely in plane $\Pi_1$, it is perpendicular to the normal vector $\vec{N}_1(A_1, B_1, C_1)$:
\begin{equation}
    A_1 l + B_1 m + C_1 n = 0. \label{eq:p39_orth1}
\end{equation}
Similarly, since $L$ lies in plane $\Pi_2$, it is perpendicular to the normal vector $\vec{N}_2(A_2, B_2, C_2)$:
\begin{equation}
    A_2 l + B_2 m + C_2 n = 0. \label{eq:p39_orth2}
\end{equation}
Solving equations \eqref{eq:p39_orth1} and \eqref{eq:p39_orth2} by the rule of cross-multiplication:
\begin{equation}
    \boxed{\frac{l}{B_1 C_2 - B_2 C_1} = \frac{m}{C_1 A_2 - C_2 A_1} = \frac{n}{A_1 B_2 - A_2 B_1}}
\end{equation}
This corresponds to the cross product of the normal vectors: $\vec{b} = \vec{N}_1 \times \vec{N}_2$.

\subsubsection{Step 2: Finding a Point on the Line}
The line $L$ intersects at least one of the coordinate planes (unless it is parallel to it).
Assuming $A_1 B_2 - A_2 B_1 \neq 0$, the line meets the $xy$-plane ($z = 0$) at some point $(x_0, y_0, 0)$.
Setting $z = 0$ in equations \eqref{eq:p39_pi1} and \eqref{eq:p39_pi2}:
\begin{align}
    A_1 x_0 + B_1 y_0 + D_1 &= 0, \\
    A_2 x_0 + B_2 y_0 + D_2 &= 0.
\end{align}
Solving this $2 \times 2$ system via Cramer's Rule:
\begin{equation}
    x_0 = \frac{B_1 D_2 - B_2 D_1}{A_1 B_2 - A_2 B_1}, \quad y_0 = \frac{D_1 A_2 - D_2 A_1}{A_1 B_2 - A_2 B_1}, \quad z_0 = 0.
\end{equation}
Thus, $(x_0, y_0, 0)$ is a specific point on the line.

\subsubsection{Step 3: Writing the Symmetrical Equation}
Combining the point $(x_0, y_0, z_0)$ and the direction ratios $(l, m, n)$:
\begin{equation}
    \boxed{\frac{x - x_0}{B_1 C_2 - B_2 C_1} = \frac{y - y_0}{C_1 A_2 - C_2 A_1} = \frac{z - 0}{A_1 B_2 - A_2 B_1}}
\end{equation}

% =========================================================================
% PAGE 40
% =========================================================================
\section{Page 40: Angle Between a Line and a Plane}
\label{sec:page40}

\begin{conceptbox}[Manuscript Overview: Page 40]
Page 40 derives the sine formula for the angle between a straight line and a plane and deduces the conditions for parallelism and perpendicularity.
\end{conceptbox}

\subsection{Derivation of the Sine Formula}
\begin{theorembox}[Angle Between a Line and a Plane]
Let $\theta$ be the angle between the straight line:
\begin{equation}
    L: \frac{x - x_1}{l} = \frac{y - y_1}{m} = \frac{z - z_1}{n}
\end{equation}
and the plane:
\begin{equation}
    \Pi: Ax + By + Cz + D = 0.
\end{equation}
Then:
\begin{equation}
    \boxed{\sin\theta = \frac{Al + Bm + Cn}{\sqrt{A^2 + B^2 + C^2}\,\sqrt{l^2 + m^2 + n^2}}}
\end{equation}
\end{theorembox}

\begin{proof}
By definition, the angle $\theta$ between a line and a plane is the angle between the line $L$ and its orthogonal projection on the plane $\Pi$.
The normal $\vec{N}$ to the plane has direction ratios $(A, B, C)$.
The angle $\phi$ between the line $L$ (direction ratios $(l, m, n)$) and the normal $\vec{N}$ is:
\begin{equation}
    \phi = 90^\circ - \theta.
\end{equation}

\begin{center}
\begin{tikzpicture}[scale=1.1, >=Stealth]
    % Plane
    \draw[thick, primary] (-3,0) -- (3,0) node[right] {Plane $\Pi$};
    \fill[primary!5] (-3,-0.1) rectangle (3,0);

    % Normal
    \draw[->, very thick, accent] (0,0) -- (0,2.5) node[above] {Normal $\vec{N}(A,B,C)$};

    % Line L
    \draw[->, very thick, secondary] (-2.5,-1) -- (2.5,1.8) node[right] {Line $L(l,m,n)$};

    % Angles
    \draw[->, thick] (1,0) arc (0:36:1) node[midway, right] {$\theta$};
    \draw[->, thick, accent] (0,1.2) arc (90:36:1.2) node[midway, above right] {$\phi = 90^\circ - \theta$};

    \node[below] at (0,0) {$O$};
\end{tikzpicture}
\end{center}

Applying the cosine formula between two directed vectors (Page 22):
\begin{equation}
    \cos\phi = \frac{Al + Bm + Cn}{\sqrt{A^2 + B^2 + C^2}\,\sqrt{l^2 + m^2 + n^2}}.
\end{equation}
Since $\cos\phi = \cos(90^\circ - \theta) = \sin\theta$:
\begin{equation}
    \boxed{\sin\theta = \frac{Al + Bm + Cn}{\sqrt{A^2 + B^2 + C^2}\,\sqrt{l^2 + m^2 + n^2}}}
\end{equation}
\end{proof}

\subsection{Conditions of Parallelism and Perpendicularity}
\begin{itemize}
    \item \textbf{Line Parallel to Plane}: The line is parallel to the plane when $\theta = 0 \iff \sin\theta = 0$:
    \begin{equation}
        \boxed{Al + Bm + Cn = 0}
    \end{equation}
    (The line is perpendicular to the normal).
    \item \textbf{Line Perpendicular to Plane}: The line is perpendicular to the plane when $\theta = 90^\circ \iff \phi = 0$:
    \begin{equation}
        \boxed{\frac{A}{l} = \frac{B}{m} = \frac{C}{n}}
    \end{equation}
    (The line is parallel to the normal).
\end{itemize}

% =========================================================================
% PAGE 41
% =========================================================================
\section{Page 41: Intersection of a Line and a Plane, Foot of Perpendicular, and Image}
\label{sec:page41}

\begin{conceptbox}[Manuscript Overview: Page 41]
Page 41 solves for:
\begin{enumerate}[noitemsep]
    \item The point of intersection between a straight line and a plane.
    \item The foot of the perpendicular dropped from an external point to a plane.
    \item The optical image (reflection) of a point in a plane.
\end{enumerate}
\end{conceptbox}

\subsection{Point of Intersection of a Line and a Plane}
Let the line be $L: \frac{x - x_1}{l} = \frac{y - y_1}{m} = \frac{z - z_1}{n} = r$, and let the plane be $\Pi: Ax + By + Cz + D = 0$.
The general point on the line is $P(r) = (x_1 + lr, \; y_1 + mr, \; z_1 + nr)$.
Substitute $P(r)$ into the plane equation:
\begin{align}
    A(x_1 + lr) + B(y_1 + mr) + C(z_1 + nr) + D &= 0 \nonumber\\
    (Ax_1 + By_1 + Cz_1 + D) + r(Al + Bm + Cn) &= 0 \nonumber\\
    \implies r &= -\frac{Ax_1 + By_1 + Cz_1 + D}{Al + Bm + Cn}.
\end{align}
Substituting this value of $r$ back into $P(r)$ yields the unique point of intersection (provided $Al + Bm + Cn \neq 0$).

\subsection{Foot of Perpendicular and Image of a Point}
\begin{theorembox}[Foot of Perpendicular and Reflection Formula]
Let $P(x_0, y_0, z_0)$ be a point and $\Pi: Ax + By + Cz + D = 0$ be a plane.
\begin{enumerate}
    \item The foot of the perpendicular $N(x, y, z)$ from $P$ to $\Pi$ is given by:
    \begin{equation}
        \boxed{\frac{x - x_0}{A} = \frac{y - y_0}{B} = \frac{z - z_0}{C} = -\frac{Ax_0 + By_0 + Cz_0 + D}{A^2 + B^2 + C^2}}
    \end{equation}
    \item The image (reflection) $P'(x', y', z')$ of $P$ in the plane $\Pi$ is given by:
    \begin{equation}
        \boxed{\frac{x' - x_0}{A} = \frac{y' - y_0}{B} = \frac{z' - z_0}{C} = -2\,\frac{Ax_0 + By_0 + Cz_0 + D}{A^2 + B^2 + C^2}}
    \end{equation}
\end{enumerate}
\end{theorembox}

\begin{proof}
The normal line passing through $P(x_0, y_0, z_0)$ has direction ratios $(A, B, C)$. Its equations are:
\begin{equation}
    \frac{x - x_0}{A} = \frac{y - y_0}{B} = \frac{z - z_0}{C} = k.
\end{equation}
Any point on this normal line is $(x_0 + Ak, \; y_0 + Bk, \; z_0 + Ck)$.
For the foot of the perpendicular $N$, this point lies on the plane:
\begin{align}
    A(x_0 + Ak) + B(y_0 + Bk) + C(z_0 + Ck) + D &= 0 \nonumber\\
    (A^2 + B^2 + C^2)k + (Ax_0 + By_0 + Cz_0 + D) &= 0 \nonumber\\
    \implies k &= -\frac{Ax_0 + By_0 + Cz_0 + D}{A^2 + B^2 + C^2}.
\end{align}
For the image point $P'$, the foot $N$ is the midpoint of $PP'$:
\begin{equation}
    N = \frac{P + P'}{2} \implies P' = 2N - P.
\end{equation}
Thus, the parameter value for $P'$ is simply $2k$.
\end{proof}

% =========================================================================
% PAGE 42
% =========================================================================
\section{Page 42: Condition of Coplanarity of Two Straight Lines}
\label{sec:page42}

\begin{conceptbox}[Manuscript Overview: Page 42]
Page 42 derives the celebrated determinant condition for two lines in space to lie in the same plane (to be coplanar) and establishes the equation of the plane containing them.
\end{conceptbox}

\subsection{Derivation of the Coplanarity Condition}
\begin{theorembox}[Coplanarity Determinant Condition]
Two straight lines:
\begin{align}
    L_1 &: \frac{x - x_1}{l_1} = \frac{y - y_1}{m_1} = \frac{z - z_1}{n_1}, \label{eq:p42_L1} \\
    L_2 &: \frac{x - x_2}{l_2} = \frac{y - y_2}{m_2} = \frac{z - z_2}{n_2}, \label{eq:p42_L2}
\end{align}
are \textbf{coplanar} if and only if:
\begin{equation}
    \boxed{\begin{vmatrix}
    x_2 - x_1 & y_2 - y_1 & z_2 - z_1 \\
    l_1 & m_1 & n_1 \\
    l_2 & m_2 & n_2
    \end{vmatrix} = 0}
\end{equation}
\end{theorembox}

\begin{proof}
Any plane passing through the line $L_1$ must pass through the point $(x_1, y_1, z_1)$ and have its normal $\vec{N}(A, B, C)$ perpendicular to $L_1$:
\begin{align}
    A(x - x_1) + B(y - y_1) + C(z - z_1) &= 0, \label{eq:p42_plane} \\
    Al_1 + Bm_1 + Cn_1 &= 0. \label{eq:p42_orth1}
\end{align}
If this plane also contains the line $L_2$:
\begin{enumerate}
    \item The point $(x_2, y_2, z_2)$ on $L_2$ must lie on the plane:
    \begin{equation}
        A(x_2 - x_1) + B(y_2 - y_1) + C(z_2 - z_1) = 0. \label{eq:p42_pt2}
    \end{equation}
    \item The normal $\vec{N}(A, B, C)$ must be perpendicular to $L_2$:
    \begin{equation}
        Al_2 + Bm_2 + Cn_2 = 0. \label{eq:p42_orth2}
    \end{equation}
\end{enumerate}
Equations \eqref{eq:p42_pt2}, \eqref{eq:p42_orth1}, and \eqref{eq:p42_orth2} form a homogeneous linear system in $A, B, C$:
\begin{equation}
    \begin{pmatrix}
    x_2 - x_1 & y_2 - y_1 & z_2 - z_1 \\
    l_1 & m_1 & n_1 \\
    l_2 & m_2 & n_2
    \end{pmatrix}
    \begin{pmatrix} A \\ B \\ C \end{pmatrix} = \begin{pmatrix} 0 \\ 0 \\ 0 \end{pmatrix}.
\end{equation}
For a non-trivial plane to exist, $(A, B, C) \neq (0, 0, 0)$, so the determinant of coefficients must vanish:
\begin{equation}
    \begin{vmatrix}
    x_2 - x_1 & y_2 - y_1 & z_2 - z_1 \\
    l_1 & m_1 & n_1 \\
    l_2 & m_2 & n_2
    \end{vmatrix} = 0.
\end{equation}
\end{proof}

\subsection{Equation of the Plane Containing Two Coplanar Lines}
Eliminating $A, B, C$ between equations \eqref{eq:p42_plane}, \eqref{eq:p42_orth1}, and \eqref{eq:p42_orth2}:
\begin{equation}
    \boxed{\begin{vmatrix}
    x - x_1 & y - y_1 & z - z_1 \\
    l_1 & m_1 & n_1 \\
    l_2 & m_2 & n_2
    \end{vmatrix} = 0}
\end{equation}

% =========================================================================
% PAGE 43
% =========================================================================
\section{Page 43: Condition for a Line to Lie Wholly in a Plane}
\label{sec:page43}

\begin{conceptbox}[Manuscript Overview: Page 43]
Page 43 analyzes the algebraic conditions ensuring that a given straight line lies completely within a specified plane.
\end{conceptbox}

\begin{theorembox}[Line Lying Wholly in a Plane]
The straight line:
\begin{equation}
    L: \frac{x - x_1}{l} = \frac{y - y_1}{m} = \frac{z - z_1}{n}
\end{equation}
lies wholly and completely in the plane:
\begin{equation}
    \Pi: Ax + By + Cz + D = 0
\end{equation}
if and only if the following two conditions are simultaneously satisfied:
\begin{align}
    \text{Condition 1} &: \quad \boxed{Ax_1 + By_1 + Cz_1 + D = 0} \quad \text{(Point lies in plane)}, \\
    \text{Condition 2} &: \quad \boxed{Al + Bm + Cn = 0} \quad \text{(Line perpendicular to normal)}.
\end{align}
\end{theorembox}

\begin{proof}
Any point on line $L$ has coordinates $(x_1 + lr, \; y_1 + mr, \; z_1 + nr)$, where $r \in \mathbb{R}$.
For the line to lie entirely in the plane, this point must satisfy the equation of the plane for \emph{every} value of $r$:
\begin{align}
    A(x_1 + lr) + B(y_1 + mr) + C(z_1 + nr) + D &\equiv 0 \nonumber\\
    (Ax_1 + By_1 + Cz_1 + D) + r(Al + Bm + Cn) &\equiv 0 \quad \forall r \in \mathbb{R}.
\end{align}
A linear polynomial in $r$ is identically zero if and only if both its constant term and the coefficient of $r$ vanish independently:
\begin{align}
    Ax_1 + By_1 + Cz_1 + D &= 0, \\
    Al + Bm + Cn &= 0.
\end{align}
\end{proof}

% =========================================================================
% PAGE 44
% =========================================================================
\section{Page 44: Plane Containing a Given Line and Parallel to Another Line}
\label{sec:page44}

\begin{conceptbox}[Manuscript Overview: Page 44]
Page 44 establishes the construction of a plane that contains a given straight line and is oriented parallel to a second straight line in space.
\end{conceptbox}

\subsection{Construction and Derivation}
Let the two lines be:
\begin{align}
    L_1 &: \frac{x - x_1}{l_1} = \frac{y - y_1}{m_1} = \frac{z - z_1}{n_1}, \\
    L_2 &: \frac{x - x_2}{l_2} = \frac{y - y_2}{m_2} = \frac{z - z_2}{n_2}.
\end{align}
We seek the plane $\Pi$ containing $L_1$ and parallel to $L_2$.
\begin{enumerate}
    \item Any plane containing $L_1$ passes through $(x_1, y_1, z_1)$:
    \begin{equation}
        A(x - x_1) + B(y - y_1) + C(z - z_1) = 0 \label{eq:p44_pl}
    \end{equation}
    and its normal $\vec{N}(A, B, C)$ is perpendicular to $L_1$:
    \begin{equation}
        Al_1 + Bm_1 + Cn_1 = 0. \label{eq:p44_c1}
    \end{equation}
    \item Since the plane is parallel to $L_2$, its normal $\vec{N}$ is also perpendicular to $L_2$:
    \begin{equation}
        Al_2 + Bm_2 + Cn_2 = 0. \label{eq:p44_c2}
    \end{equation}
\end{enumerate}
Eliminating $A, B, C$ between equations \eqref{eq:p44_pl}, \eqref{eq:p44_c1}, and \eqref{eq:p44_c2}:
\begin{equation}
    \boxed{\begin{vmatrix}
    x - x_1 & y - y_1 & z - z_1 \\
    l_1 & m_1 & n_1 \\
    l_2 & m_2 & n_2
    \end{vmatrix} = 0}
\end{equation}

% =========================================================================
% PAGE 45
% =========================================================================
\section{Page 45: Line Intersecting Two Given Lines and Satisfying a Geometric Condition}
\label{sec:page45}

\begin{conceptbox}[Manuscript Overview: Page 45]
Page 45 details the method of finding a transversal line that intersects two given skew lines and either passes through a fixed point or is parallel to a given direction.
\end{conceptbox}

\subsection{Transversal Passing Through a Given Point}
Let the two given skew lines be $L_1$ and $L_2$, and let the fixed point be $P(\alpha, \beta, \gamma)$.
\begin{enumerate}
    \item \textbf{Plane 1 ($\Pi_1$)}: Determine the plane containing $L_1$ and the point $P$.
    \item \textbf{Plane 2 ($\Pi_2$)}: Determine the plane containing $L_2$ and the point $P$.
    \item \textbf{Line of Intersection}: The required transversal line is the line of intersection of planes $\Pi_1$ and $\Pi_2$:
    \begin{equation}
        \Pi_1 = 0 = \Pi_2.
    \end{equation}
    Since this line lies in $\Pi_1$, it intersects $L_1$. Since it lies in $\Pi_2$, it intersects $L_2$. Since both planes pass through $P$, the line passes through $P$.
\end{enumerate}

\subsection{Transversal Parallel to a Given Line}
If the line must be parallel to a given direction $(l, m, n)$:
\begin{enumerate}
    \item Find plane $\Pi_1$ containing $L_1$ and parallel to direction $(l, m, n)$.
    \item Find plane $\Pi_2$ containing $L_2$ and parallel to direction $(l, m, n)$.
    \item The line of intersection $\Pi_1 = 0 = \Pi_2$ is the required line.
\end{enumerate}

% =========================================================================
% PAGE 46
% =========================================================================
\section{Page 46: Skew Lines and the Concept of Shortest Distance}
\label{sec:page46}

\begin{conceptbox}[Manuscript Overview: Page 46]
Page 46 introduces \textbf{skew lines}, establishes that there exists a unique common perpendicular between them, and defines the \textbf{Shortest Distance} (S.D.).
\end{conceptbox}

\subsection{Definition of Skew Lines}
\begin{conceptbox}[Skew Lines]
Two straight lines in three-dimensional space which are neither parallel nor intersecting are called \textbf{skew lines} (or non-coplanar lines).
\end{conceptbox}

\subsection{The Common Perpendicular}
\begin{theorembox}[Existence and Uniqueness of the Common Perpendicular]
For any two non-parallel, non-intersecting lines $L_1$ and $L_2$ in space, there exists a unique straight line segment $GH$ which intersects both lines perpendicularly. The length of this segment $GH$ is the \textbf{shortest distance} between $L_1$ and $L_2$.
\end{theorembox}

\begin{center}
\begin{tikzpicture}[scale=1.2, >=Stealth, 3d view={120}{20}]
    % Two parallel planes enclosing skew lines
    \fill[primary!5, opacity=0.8] (-3,-2,2) -- (3,-2,2) -- (3,2,2) -- (-3,2,2) -- cycle;
    \fill[secondary!5, opacity=0.8] (-3,-2,-1) -- (3,-2,-1) -- (3,2,-1) -- (-3,2,-1) -- cycle;

    \draw[dashed, gray] (-3,-2,2) -- (3,-2,2) -- (3,2,2) -- (-3,2,2) -- cycle;
    \draw[dashed, gray] (-3,-2,-1) -- (3,-2,-1) -- (3,2,-1) -- (-3,2,-1) -- cycle;

    % Line L1 in top plane
    \draw[very thick, primary] (-2.5, -1.2, 2) -- (2.5, 1.2, 2) node[right] {$L_1$};
    \coordinate (G) at (0,0,2);
    \fill[primary] (G) circle (2pt) node[above right] {$G$};

    % Line L2 in bottom plane
    \draw[very thick, secondary] (-2, 1.5, -1) -- (2, -1.5, -1) node[right] {$L_2$};
    \coordinate (H) at (0,0,-1);
    \fill[secondary] (H) circle (2pt) node[below right] {$H$};

    % Common perpendicular GH
    \draw[very thick, accent] (H) -- (G) node[midway, right] {$d = GH \text{ (S.D.)}$};

    % Right angle markers
    \draw[thick, accent] (0,0.3,2) -- (0.3,0.3,2) -- (0.3,0,2);
    \draw[thick, accent] (0,0.3,-1) -- (0.3,0.3,-1) -- (0.3,0,-1);
\end{tikzpicture}
\end{center}

% =========================================================================
% PAGE 47
% =========================================================================
\section{Page 47: Geometric Configuration and Direction of the Common Perpendicular}
\label{sec:page47}

\begin{conceptbox}[Manuscript Overview: Page 47]
Page 47 derives the direction cosines $(\lambda, \mu, \nu)$ of the common perpendicular to two skew lines.
\end{conceptbox}

\subsection{Direction Cosines of the Common Perpendicular}
Let the two skew lines be:
\begin{align}
    L_1 &: \frac{x - x_1}{l_1} = \frac{y - y_1}{m_1} = \frac{z - z_1}{n_1}, \\
    L_2 &: \frac{x - x_2}{l_2} = \frac{y - y_2}{m_2} = \frac{z - z_2}{n_2}.
\end{align}
Let $(\lambda, \mu, \nu)$ be the direction cosines of the line of shortest distance $GH$.
Since $GH \perp L_1$ and $GH \perp L_2$:
\begin{align}
    \lambda l_1 + \mu m_1 + \nu n_1 &= 0, \label{eq:p47_gh1} \\
    \lambda l_2 + \mu m_2 + \nu n_2 &= 0. \label{eq:p47_gh2}
\end{align}
By the method of cross-multiplication:
\begin{equation}
    \frac{\lambda}{m_1 n_2 - m_2 n_1} = \frac{\mu}{n_1 l_2 - n_2 l_1} = \frac{\nu}{l_1 m_2 - l_2 m_1} = k.
\end{equation}
Since $\lambda^2 + \mu^2 + \nu^2 = 1$:
\begin{equation}
    k = \frac{1}{\sqrt{(m_1 n_2 - m_2 n_1)^2 + (n_1 l_2 - n_2 l_1)^2 + (l_1 m_2 - l_2 m_1)^2}} = \frac{1}{\sin\theta}
\end{equation}
where $\theta$ is the angle between $L_1$ and $L_2$ (Page 23).
Thus:
\begin{equation}
    \boxed{\lambda = \frac{m_1 n_2 - m_2 n_1}{\sin\theta}, \quad \mu = \frac{n_1 l_2 - n_2 l_1}{\sin\theta}, \quad \nu = \frac{l_1 m_2 - l_2 m_1}{\sin\theta}}
\end{equation}

% =========================================================================
% PAGE 48
% =========================================================================
\section{Page 48: Analytical Derivation of the Magnitude of Shortest Distance}
\label{sec:page48}

\begin{conceptbox}[Manuscript Overview: Page 48]
Page 48 delivers the complete analytical derivation of the formula for the magnitude of the shortest distance between two skew lines.
\end{conceptbox}

\begin{theorembox}[Magnitude of Shortest Distance (S.D.)]
The shortest distance $d$ between the skew lines $L_1$ and $L_2$ is:
\begin{equation}
    \boxed{d = \frac{\left|\begin{matrix}
    x_2 - x_1 & y_2 - y_1 & z_2 - z_1 \\
    l_1 & m_1 & n_1 \\
    l_2 & m_2 & n_2
    \end{matrix}\right|}{\sqrt{(m_1 n_2 - m_2 n_1)^2 + (n_1 l_2 - n_2 l_1)^2 + (l_1 m_2 - l_2 m_1)^2}}}
\end{equation}
\end{theorembox}

\begin{proof}
Let $P(x_1, y_1, z_1)$ be a point on $L_1$ and $Q(x_2, y_2, z_2)$ be a point on $L_2$.
The common perpendicular segment $GH$ is perpendicular to both lines.
Hence, the length $d = GH$ is equal to the orthogonal projection of the displacement segment $PQ$ onto the directed line $GH$ whose direction cosines are $(\lambda, \mu, \nu)$ (Page 20):
\begin{equation}
    d = |\text{Proj}_{GH}(\vec{PQ})| = |\lambda(x_2 - x_1) + \mu(y_2 - y_1) + \nu(z_2 - z_1)|.
\end{equation}
Substituting the expressions for $\lambda, \mu, \nu$ from Page 47:
\begin{align}
    d &= \left| (x_2 - x_1)\frac{m_1 n_2 - m_2 n_1}{\sin\theta} + (y_2 - y_1)\frac{n_1 l_2 - n_2 l_1}{\sin\theta} + (z_2 - z_1)\frac{l_1 m_2 - l_2 m_1}{\sin\theta} \right| \nonumber\\
    &= \frac{1}{\sin\theta} \left| (x_2 - x_1)(m_1 n_2 - m_2 n_1) + (y_2 - y_1)(n_1 l_2 - n_2 l_1) + (z_2 - z_1)(l_1 m_2 - l_2 m_1) \right|.
\end{align}
Recognizing the numerator as the expansion of a $3 \times 3$ determinant:
\begin{equation}
    d = \frac{\left|\begin{matrix}
    x_2 - x_1 & y_2 - y_1 & z_2 - z_1 \\
    l_1 & m_1 & n_1 \\
    l_2 & m_2 & n_2
    \end{matrix}\right|}{\sqrt{(m_1 n_2 - m_2 n_1)^2 + (n_1 l_2 - n_2 l_1)^2 + (l_1 m_2 - l_2 m_1)^2}}.
\end{equation}
\end{proof}

\begin{notebox}[Intersection as a Degenerate Case ($d = 0$)]
If the two lines intersect, the shortest distance between them is zero:
\begin{equation}
    d = 0 \iff \begin{vmatrix}
    x_2 - x_1 & y_2 - y_1 & z_2 - z_1 \\
    l_1 & m_1 & n_1 \\
    l_2 & m_2 & n_2
    \end{vmatrix} = 0.
\end{equation}
This perfectly reproduces the coplanarity condition established on Page 42!
\end{notebox}

% =========================================================================
% PAGE 49
% =========================================================================
\section{Page 49: Equations of the Line of Shortest Distance}
\label{sec:page49}

\begin{conceptbox}[Manuscript Overview: Page 49]
Page 49 constructs the equations of the actual straight line containing the common perpendicular $GH$, expressed as the intersection of two planes.
\end{conceptbox}

\begin{theorembox}[Equations of the Line of Shortest Distance]
The line of shortest distance between $L_1$ and $L_2$ is represented by the pair of planes:
\begin{equation}
    \boxed{\begin{vmatrix}
    x - x_1 & y - y_1 & z - z_1 \\
    l_1 & m_1 & n_1 \\
    \lambda & \mu & \nu
    \end{vmatrix} = 0 \quad \text{and} \quad
    \begin{vmatrix}
    x - x_2 & y - y_2 & z - z_2 \\
    l_2 & m_2 & n_2 \\
    \lambda & \mu & \nu
    \end{vmatrix} = 0}
\end{equation}
where $(\lambda, \mu, \nu)$ are the direction ratios of the common perpendicular.
\end{theorembox}

\begin{proof}
The line of shortest distance $GH$ intersects $L_1$ at $G$ and is perpendicular to $L_1$.
Thus, $GH$ and $L_1$ are coplanar lines intersecting at $G$.
The plane containing both $L_1$ and $GH$ passes through $(x_1, y_1, z_1)$ and has directions $(l_1, m_1, n_1)$ and $(\lambda, \mu, \nu)$. Its equation is:
\begin{equation}
    \Pi_1: \begin{vmatrix}
    x - x_1 & y - y_1 & z - z_1 \\
    l_1 & m_1 & n_1 \\
    \lambda & \mu & \nu
    \end{vmatrix} = 0.
\end{equation}
Similarly, $GH$ intersects $L_2$ at $H$ and is perpendicular to $L_2$.
The plane containing both $L_2$ and $GH$ passes through $(x_2, y_2, z_2)$ and has directions $(l_2, m_2, n_2)$ and $(\lambda, \mu, \nu)$. Its equation is:
\begin{equation}
    \Pi_2: \begin{vmatrix}
    x - x_2 & y - y_2 & z - z_2 \\
    l_2 & m_2 & n_2 \\
    \lambda & \mu & \nu
    \end{vmatrix} = 0.
\end{equation}
Since the line $GH$ lies in both planes $\Pi_1$ and $\Pi_2$, it is precisely their line of intersection $\Pi_1 = 0 = \Pi_2$.
\end{proof}

% =========================================================================
% PAGE 50
% =========================================================================
\section{Page 50: Comprehensive Solved Examination Problem on Skew Lines}
\label{sec:page50}

\begin{conceptbox}[Manuscript Overview: Page 50]
Page 50 works through an extensive university examination problem, calculating the shortest distance between two given skew lines.
\end{conceptbox}

\begin{examplebox}[Examination Problem: Shortest Distance Calculation]
Find the magnitude of the shortest distance between the lines:
\begin{align}
    L_1 &: \frac{x - 3}{3} = \frac{y - 8}{-1} = \frac{z - 3}{1}, \label{eq:ex_L1} \\
    L_2 &: \frac{x + 3}{-3} = \frac{y + 7}{2} = \frac{z - 6}{4}. \label{eq:ex_L2}
\end{align}
\end{examplebox}

\begin{proof}[Step-by-Step Analytical Solution]
\textbf{Step 1: Identify Points and Direction Ratios}:
\begin{align}
    \text{From } L_1 &: \quad P_1 = (3, 8, 3), \quad (l_1, m_1, n_1) = (3, -1, 1), \\
    \text{From } L_2 &: \quad P_2 = (-3, -7, 6), \quad (l_2, m_2, n_2) = (-3, 2, 4).
\end{align}

\textbf{Step 2: Direction Ratios of the Common Perpendicular}:
Let the common perpendicular have direction ratios $(\lambda, \mu, \nu)$:
\begin{align}
    3\lambda - \mu + \nu &= 0, \\
    -3\lambda + 2\mu + 4\nu &= 0.
\end{align}
Applying the cross-multiplication rule:
\begin{align}
    \frac{\lambda}{(-1)(4) - (1)(2)} &= \frac{\mu}{(1)(-3) - (3)(4)} = \frac{\nu}{(3)(2) - (-1)(-3)} \nonumber\\
    \frac{\lambda}{-4 - 2} &= \frac{\mu}{-3 - 12} = \frac{\nu}{6 - 3} \nonumber\\
    \frac{\lambda}{-6} &= \frac{\mu}{-15} = \frac{\nu}{3}.
\end{align}
Dividing by $-3$:
\begin{equation}
    \boxed{\frac{\lambda}{2} = \frac{\mu}{5} = \frac{\nu}{-1}}
\end{equation}
Thus, a set of direction ratios for the common perpendicular is $(a, b, c) = (2, 5, -1)$.

\textbf{Step 3: Normalize to Direction Cosines}:
\begin{align}
    \sqrt{a^2 + b^2 + c^2} &= \sqrt{2^2 + 5^2 + (-1)^2} = \sqrt{4 + 25 + 1} = \sqrt{30}. \\
    (\lambda, \mu, \nu) &= \left(\frac{2}{\sqrt{30}}, \; \frac{5}{\sqrt{30}}, \; \frac{-1}{\sqrt{30}}\right).
\end{align}

\textbf{Step 4: Vector Connecting Fixed Points}:
\begin{equation}
    \vec{P_1 P_2} = (-3 - 3)\,\hat{i} + (-7 - 8)\,\hat{j} + (6 - 3)\,\hat{k} = -6\,\hat{i} - 15\,\hat{j} + 3\,\hat{k}.
\end{equation}

\textbf{Step 5: Compute the Shortest Distance}:
\begin{align}
    d &= |\lambda(x_2 - x_1) + \mu(y_2 - y_1) + \nu(z_2 - z_1)| \nonumber\\
    &= \left| \frac{2}{\sqrt{30}}(-6) + \frac{5}{\sqrt{30}}(-15) + \frac{-1}{\sqrt{30}}(3) \right| \nonumber\\
    &= \frac{|-12 - 75 - 3|}{\sqrt{30}} = \frac{|-90|}{\sqrt{30}} = \frac{90}{\sqrt{30}}.
\end{align}
Rationalizing the denominator:
\begin{equation}
    \boxed{d = \frac{90\sqrt{30}}{30} = 3\sqrt{30}}
\end{equation}
\end{proof}

% =========================================================================
% PAGE 51
% =========================================================================
\section{Page 51: Equations of the Line of S.D. and Feet of Common Perpendicular}
\label{sec:page51}

\begin{conceptbox}[Manuscript Overview: Page 51]
Page 51 completes the examination problem by:
\begin{enumerate}[noitemsep]
    \item Finding the Cartesian equations of the line of shortest distance.
    \item Determining the exact spatial coordinates of the points of intersection (feet) $G$ and $H$.
\end{enumerate}
\end{conceptbox}

\subsection{Equations of the Line of Shortest Distance}
Applying the two-plane representation from Page 49 with $(\lambda, \mu, \nu) \sim (2, 5, -1)$:
\begin{enumerate}
    \item \textbf{Plane 1 (Containing $L_1$ and $GH$)}:
    \begin{equation}
        \begin{vmatrix}
        x - 3 & y - 8 & z - 3 \\
        3 & -1 & 1 \\
        2 & 5 & -1
        \end{vmatrix} = 0.
    \end{equation}
    Expanding along row 1:
    \begin{align}
        (x - 3)[(-1)(-1) - (1)(5)] - (y - 8)[(3)(-1) - (1)(2)] + (z - 3)[(3)(5) - (-1)(2)] &= 0 \nonumber\\
        (x - 3)(1 - 5) - (y - 8)(-3 - 2) + (z - 3)(15 + 2) &= 0 \nonumber\\
        -4(x - 3) + 5(y - 8) + 17(z - 3) &= 0 \nonumber\\
        -4x + 12 + 5y - 40 + 17z - 51 &= 0 \nonumber\\
        \Aboxed{4x - 5y - 17z + 79 &= 0} \label{eq:p51_pl1}
    \end{align}
    \item \textbf{Plane 2 (Containing $L_2$ and $GH$)}:
    \begin{equation}
        \begin{vmatrix}
        x + 3 & y + 7 & z - 6 \\
        -3 & 2 & 4 \\
        2 & 5 & -1
        \end{vmatrix} = 0.
    \end{equation}
    Expanding along row 1:
    \begin{align}
        (x + 3)[(2)(-1) - (4)(5)] - (y + 7)[(-3)(-1) - (4)(2)] + (z - 6)[(-3)(5) - (2)(2)] &= 0 \nonumber\\
        (x + 3)(-2 - 20) - (y + 7)(3 - 8) + (z - 6)(-15 - 4) &= 0 \nonumber\\
        -22(x + 3) + 5(y + 7) - 19(z - 6) &= 0 \nonumber\\
        -22x - 66 + 5y + 35 - 19z + 114 &= 0 \nonumber\\
        \Aboxed{22x - 5y + 19z - 83 &= 0} \label{eq:p51_pl2}
    \end{align}
\end{enumerate}
The equations of the line of shortest distance are given by \eqref{eq:p51_pl1} and \eqref{eq:p51_pl2} simultaneously.

\subsection{Coordinates of the Feet $G$ and $H$}
Let $G$ be a point on $L_1$ and $H$ be a point on $L_2$:
\begin{align}
    G &= (3 + 3r, \; 8 - r, \; 3 + r), \\
    H &= (-3 - 3s, \; -7 + 2s, \; 6 + 4s).
\end{align}
The direction ratios of the segment $GH$ are:
\begin{equation}
    \vec{GH} = H - G = \bigl(-6 - 3s - 3r, \; -15 + 2s + r, \; 3 + 4s - r\bigr).
\end{equation}
Since $GH$ is along the common perpendicular with DRs $(2, 5, -1)$:
\begin{equation}
    \frac{-6 - 3s - 3r}{2} = \frac{-15 + 2s + r}{5} = \frac{3 + 4s - r}{-1} = k.
\end{equation}
Solving this linear system yields $r = 0$ and $s = 0$!
Hence:
\begin{equation}
    \boxed{G = (3, 8, 3), \quad H = (-3, -7, 6)}
\end{equation}
\textbf{Verification}:
\begin{equation}
    GH = \sqrt{(-3-3)^2 + (-7-8)^2 + (6-3)^2} = \sqrt{(-6)^2 + (-15)^2 + 3^2} = \sqrt{36 + 225 + 9} = \sqrt{270} = 3\sqrt{30}.
\end{equation}
The distance matches the formula from Page 50 precisely!

% =========================================================================
% PAGE 52
% =========================================================================
\section{Page 52: Distance Between Parallel Lines and Chapter Summary}
\label{sec:page52}

\begin{conceptbox}[Manuscript Overview: Page 52]
Page 52 concludes the analytical treatment of lines by considering the case of \textbf{parallel lines} (where $\vec{b}_1 \times \vec{b}_2 = \vec{0}$) and provides a comparative summary.
\end{conceptbox}

\subsection{Shortest Distance Between Parallel Lines}
If two lines $L_1$ and $L_2$ are parallel, their direction ratios are proportional:
\begin{align}
    L_1 &: \frac{x - x_1}{l} = \frac{y - y_1}{m} = \frac{z - z_1}{n}, \\
    L_2 &: \frac{x - x_2}{l} = \frac{y - y_2}{m} = \frac{z - z_2}{n}.
\end{align}
Here, the lines lie in a common plane, and the cross product of their directions vanishes.
The perpendicular distance $d$ between them is obtained by taking the projection of $\vec{P_1 P_2}$ perpendicular to the direction vector $\hat{u}(l, m, n)$:

\begin{theorembox}[Distance Between Parallel Lines]
The perpendicular distance between two parallel lines with unit direction vector $\hat{u} = l\,\hat{i} + m\,\hat{j} + n\,\hat{k}$ is:
\begin{equation}
    \boxed{d = |\vec{P_1 P_2} \times \hat{u}| = \frac{\sqrt{\begin{vmatrix} y_2-y_1 & z_2-z_1 \\ m & n \end{vmatrix}^2 + \begin{vmatrix} z_2-z_1 & x_2-x_1 \\ n & l \end{vmatrix}^2 + \begin{vmatrix} x_2-x_1 & y_2-y_1 \\ l & m \end{vmatrix}^2}}{\sqrt{l^2 + m^2 + n^2}}}
\end{equation}
\end{theorembox}

\subsection{Summary of Spatial Line Geometry}
\begin{center}
\begin{tabular}{|l|l|}
\hline
\textbf{Geometric Entity} & \textbf{Standard Analytical Expression} \\
\hline
Symmetrical line & $\frac{x-x_1}{l} = \frac{y-y_1}{m} = \frac{z-z_1}{n} = r$ \\
General point on line & $(x_1 + lr, \; y_1 + mr, \; z_1 + nr)$ \\
Angle with plane & $\sin\theta = \frac{Al+Bm+Cn}{\sqrt{\sum A^2}\sqrt{\sum l^2}}$ \\
Coplanarity condition & $\det[\vec{P_2}-\vec{P_1}, \; \vec{b}_1, \; \vec{b}_2] = 0$ \\
Shortest distance (skew) & $d = \frac{|\det[\vec{P_2}-\vec{P_1}, \; \vec{b}_1, \; \vec{b}_2]|}{|\vec{b}_1 \times \vec{b}_2|}$ \\
Distance (parallel) & $d = \frac{|(\vec{P_2}-\vec{P_1}) \times \vec{b}|}{|\vec{b}|}$ \\
\hline
\end{tabular}
\end{center}
